
\section{Explicit sixth and seventh normalized scalar moments}
\label{hsm:section}

The all-degree bridge in Section~\ref{gat:section} specifies a finite
calculation at each weight.  We now carry it out at degrees six and
seven, auditing the explicit reductions supplied by the
\emph{Differential Polygamma and Lerch Identities} continuation in
\cite{dpm:reports}.  Every retained mixed derivative is a convergent
coordinate of the strict multiple zeta function.  We do not claim a
further arithmetic elimination of those coordinates.

For an admissible word $\mathbf w$, put
$Z_{\mathbf w}=Z(\mathbf w)$ and
$\Delta_{ij}Z_{\mathbf w}=(\partial_i-\partial_j)Z(\mathbf w)$.
The partial derivatives belong to a genuinely free-order holomorphic
function near the displayed indices.  They are not derivatives of
integer-index duality relations.  Retain $\alpha_j,L_d,\mathsf P_m$
and $\mathsf S_m$ exactly as defined in Section~\ref{gat:section}.

\begin{theorem}[The sixth normalized moment]
\label{hsm:sixth}
The original-coordinate finite part is
\begin{equation}\label{hsm:sixth-main}
 \boxed{\mathsf S_6=720(\alpha_6-L_5)
          +245\zeta(2)\zeta(3)+175\zeta(5)+\mathsf D_6,}
\end{equation}
where the complete convergent correction is
\begin{align}\label{hsm:D6}
 \mathsf D_6={}&120\zeta(2)\Delta_{12}Z_{2,1}
       -60\Delta_{23}Z_{2,2,1}+120\Delta_{12}Z_{3,1,1}\notag\\
 &-180\Delta_{12}Z_{4,1}-30\Delta_{12}Z_{2,3}
       +138\zeta'(5)-60\zeta(2)\zeta'(3)+30\zeta(3)\zeta'(2).
\end{align}
\end{theorem}

\begin{theorem}[The seventh normalized moment]
\label{hsm:seventh}
In the same normalization,
\begin{equation}\label{hsm:seventh-main}
 \boxed{\mathsf S_7=5040(\alpha_7-L_6)
      +\frac{3731}{4}\zeta(3)^2+\frac{918631}{240}\zeta(6)+\mathsf D_7,}
\end{equation}
with
\begin{align}\label{hsm:D7}
 \mathsf D_7={}&840\zeta(2)\Delta_{12}Z_{2,1,1}
       -840\zeta(2)\Delta_{12}Z_{3,1}+210\zeta(2)\partial_1Z_{2,2}\notag\\
 &-420\Delta_{23}Z_{2,2,1,1}+630\Delta_{23}Z_{2,3,1}
       +840\Delta_{12}Z_{3,1,1,1}\notag\\
 &+420\Delta_{12}Z_{3,1,2}
       +210(2\Delta_{13}+\Delta_{23})Z_{3,2,1}
       +210\partial_1Z_{2,2,2}\notag\\
 &+798\partial_1Z_{3,3}-1260\Delta_{12}Z_{4,1,1}
       +168\Delta_{12}Z_{2,4}+1596\Delta_{12}Z_{5,1}\notag\\
 &+\frac{651}{2}\zeta(4)\zeta'(2)+336\zeta(2)\zeta'(4)
       -378\zeta(3)\zeta'(3)-504\zeta'(6).
\end{align}
\end{theorem}

\subsection{Exact pole extraction and the marked-word ledger}

Let $h=H_n$ and $A_r=\zeta(r)-H_n^{(r)}$, the strict upper tails.
The definitions in Section~\ref{gat:section} give
\begin{align}\label{hsm:delta6}
 \delta_6={}&30A_2^2h-90A_2A_3-120\zeta(2)A_2h
       -60A_3h^2+120\zeta(2)A_3+150A_4h-138A_5,\\
 \delta_7={}&70A_2^3+105A_2^2h^2-315\zeta(2)A_2^2
       -630A_2A_3h+588A_2A_4\notag\\
 &-420\zeta(2)A_2h^2-126\zeta(4)A_2+259A_3^2
       -140A_3h^3+840\zeta(2)A_3h-280\zeta(3)A_3\notag\\
 &+525A_4h^2-651\zeta(2)A_4-966A_5h+833A_6.
       \label{hsm:delta7}
\end{align}
These are finite polynomial identities from the convergent local
series and the Gamma--Appell coefficient generator; even zeta
relations are enforced.  Apply the exact difference rule
\eqref{gat:difference}, then the upper/lower stuffle products and
marked concatenation \eqref{gat:mark}.  The resulting unsimplified
word coefficients are recorded independently in
\path{verification/appell-six-seven-exact.json}.  There are thirteen
nonzero marked coordinates at degree six and thirty-two at degree
seven before the free-order reductions below.

Only identities valid on open convergence domains are differentiated:
\[
 \zeta(a)\zeta(b)=Z(a,b)+Z(b,a)+\zeta(a+b),\qquad
 Z(s,s,s)=\frac{\zeta(s)^3-3\zeta(s)\zeta(2s)+2\zeta(3s)}6.
\]
In particular,
\begin{align*}
 (\partial_1+\partial_2)Z_{2,2}
       &=\zeta(2)\zeta'(2)-\zeta'(4),\\
 (\partial_1+\partial_2)Z_{3,3}
       &=\zeta(3)\zeta'(3)-\zeta'(6),\\
 (\partial_1+\partial_2+\partial_3)Z_{2,2,2}
       &=\frac34\zeta(4)\zeta'(2)-\zeta(2)\zeta'(4)+\zeta'(6).
\end{align*}
Differentiating the first identity in either argument also transfers
the $(3,2)$ and $(4,2)$ coordinates into $(2,3)$ and $(2,4)$.
Substitution in the generated word lists gives
\eqref{hsm:D6}--\eqref{hsm:D7} exactly.  Every index word is
admissible and every marked coordinate converges in a full neighborhood.

\subsection{The ordinary constants have finite rational certificates}

The independent higher-pole compiler gives
\begin{align}\label{hsm:J6}
 \mathsf J_6={}&75\zeta(2)\zeta(3)+60\zeta(2)Z_{2,1}
       +360Z_{2,1,1,1}+60Z_{2,1,2}+60Z_{2,2,1}\notag\\
 &-30Z_{2,3}-15Z_{3,2}-35Z_{4,1}-\frac52\zeta(5),\\
 \mathsf J_7={}&511\zeta(2)\zeta(4)+630\zeta(3)Z_{2,1}
       +420\zeta(2)Z_{2,1,1}+2520Z_{2,1,1,1,1}\notag\\
 &+420(Z_{2,1,1,2}+Z_{2,1,2,1}+Z_{2,2,1,1})
       -210(Z_{2,1,3}+Z_{2,3,1})-42Z_{2,4}\notag\\
 &-42\zeta(3)^2-105\zeta(2)Z_{3,1}
       -105(Z_{3,1,2}+Z_{3,2,1})-63Z_{3,3}\notag\\
 &-245Z_{4,1,1}-77Z_{4,2}-\frac{35}{2}Z_{5,1}
       +\frac{189}{10}\zeta(6).
       \label{hsm:J7}
\end{align}
The coefficients $75$ and $511$ combine the like terms in the
unreduced ledger; no new zeta relation is used in that combination.

Here is a reproducible finite derivation of the ordinary reductions.
At weight $w=5$ or $6$, enumerate every positive-index word of total
weight $w$ whose first entry is at least two.  Encode an entry $k$
as $0^{k-1}1$.  The series stuffle product, binary integral shuffle
product, and reversal-complement duality give exact linear equations.
Use all convergent products of weights $k+(w-k)$, with $k\ge2$
and $w-k\ge2$, and the normalization $Z_w=\zeta(w)$.
Add the shuffle-minus-stuffle relation of $(1)$ with every admissible
weight-$(w-1)$ word.  Its nonadmissible terms cancel.

The validity of this last step is specific to a single logarithmic
divergence.  The classical comparison theorem
\cite[Theorem~1, printed p.~310]{dgp:IKZ} relates the two polynomial
regularizations by an operator whose exponential generating factor
is $1+O(u^2)$.  It is therefore the identity on constants and linear
polynomials in the common divergence parameter.  Each word here has
at most one leading $1$, so the common linear product cancels with
no extra quadratic regularization term.  Equating higher powers of
the divergence parameter without that comparison would be invalid.
For example this construction gives $Z_{2,1}=\zeta(3)$ from the
single-divergence product of $(1)$ and $(2)$.

The independently generated systems have $21$ rows and $8$ columns
at weight five, and $49$ rows and $16$ columns at weight six.  Both
have full column rank over $\mathbb Q$.  The receipt includes all
rows, all right sides, a selected square minor and its inverse.
Both inverse products and every unselected row residual are checked.
The entire delivered weight-six matrix is independently regenerated.
These are representations in specified ordinary constants, without
an independence assertion about the evaluated constants themselves.

At weight five put $c=\zeta(2)\zeta(3)$ and $d=\zeta(5)$.
The required values are
\[
 Z_{2,1,1,1}=d,\quad
 Z_{2,1,2}=Z_{2,3}=\tfrac92d-2c,\quad
 Z_{2,2,1}=Z_{3,2}=3c-\tfrac{11}{2}d,\quad
 Z_{4,1}=2d-c.
\]
At weight six put $a=\zeta(3)^2$ and $b=\zeta(6)$:
\[
\begin{array}{c|l}
 Z_{2,1,1,1,1}&b\\
 Z_{2,1,1,2}=Z_{2,4}&25b/12-a\\
 Z_{2,1,2,1}=Z_{3,3}&(a-b)/2\\
 Z_{2,2,1,1}=Z_{4,2}&a-4b/3\\
 Z_{2,1,3}&a-13b/16\\
 Z_{2,3,1}=Z_{3,1,2}&53b/24-3a/2\\
 Z_{3,2,1}&3a-203b/48\\
 Z_{4,1,1}&23b/16-a\\
 Z_{5,1}&3b/4-a/2.
\end{array}
\]
Together with the already established lower-weight values and
$\zeta(2)\zeta(4)=7\zeta(6)/4$, substitution gives
\[
 \mathsf J_6=245\zeta(2)\zeta(3)+175\zeta(5),\qquad
 \mathsf J_7=\frac{3731}{4}\zeta(3)^2+\frac{918631}{240}\zeta(6).
\]

\begin{proof}[Proof of Theorems~\ref{hsm:sixth} and \ref{hsm:seventh}]
Use Theorem~\ref{gat:main} at $m=6,7$.  The local coefficient
calculation gives \eqref{hsm:delta6}--\eqref{hsm:delta7}, whose finite
differences and open-domain stuffle derivatives give the displayed
mixed corrections.  The exact higher-pole calculation and the valid
finite ordinary systems give \eqref{hsm:J6}--\eqref{hsm:J7} and their
reductions.  All convergence and fixed-coordinate endpoint operations
are already proved in the all-degree bridge.  This proves the two
identities without extrapolating numerical values or differentiating
integer-only duality.
\end{proof}

For explicit integrands put $U=\psi(x)+\gamma$.  Then
\begin{align*}
 \mathsf P_6(\psi)&=U^6+15\zeta(2)U^4+40\zeta(3)U^3
       +\frac{405}{2}\zeta(4)U^2\\
 &\quad +(120\zeta(2)\zeta(3)+144\zeta(5))U
       +40\zeta(3)^2+\frac{2745}{8}\zeta(6),\\
 \mathsf P_7(\psi)&=U^7+21\zeta(2)U^5+70\zeta(3)U^4
       +\frac{945}{2}\zeta(4)U^3\\
 &\quad +(420\zeta(2)\zeta(3)+504\zeta(5))U^2
       +\left(280\zeta(3)^2+\frac{19215}{8}\zeta(6)\right)U\\
 &\quad +945\zeta(3)\zeta(4)+504\zeta(2)\zeta(5)+720\zeta(7).
\end{align*}
Equation~\eqref{gat:inverse} now recovers the raw $Q_6,Q_7$ from
these identities and the already evaluated lower moments.  It
introduces no undefined higher moment or unspecified mixed constant.
The remaining arithmetic question is reduction of the retained
convergent directions to a chosen smaller alphabet.

The positive Cauchy form replaces the first terms in the boxes by
$-720\sum\rho_n/n^7$ and $-5040\sum\rho_n/n^8$ respectively.
Independent native quadrature and residue-defect tail summation give
approximately
\[
 \begin{aligned}
 \mathsf S_6&=20.8226131542578224267340551051333\ldots,\\
 \mathsf S_7&=123.619906884273853180167345465232\ldots.
 \end{aligned}
\]
The numerical checks retain two different head/tail configurations and
an independent positive Cauchy Mellin integral.  Their agreement is a
diagnostic, not a total interval accuracy certificate.  These scalar
moments do not settle the Gaussian $S_6$ or revised $S_8$ candidates.
